% GENERATED FILE. Edit canonical lessons, then run npm run generate:books.
% canonical-source-hash: fnv1a64:ea4e35df5ceb4870
% canonical-lessons: RU-C173-shit, RU-C173-rezat, RU-C173-varit, RU-C173-zharit, RU-C173-nalivat

\chapter{To sew, To cut, To boil, To fry, To pour}
\label{ch:ru-to-sew-to-cut-to-boil-to-fry-to-pour}
\hlchaptermodality{173}

\begin{chapteropening}
\textbf{By the end of this chapter:} \emph{I can say to sew, to cut, to boil, to fry and to pour.}

Complete the last lesson of chapter 173: I can say to sew, to cut, to boil, to fry and to pour.
\end{chapteropening}

\section[\texorpdfstring{\ru{шить} (shit)}{шить (shit)}]{\ru{шить} (shit) — to sew}

\label{lesson:RU-C173-shit}



\begin{quote}
Before the new one: say the Russian for a business trip, then the Russian for a contract.
\end{quote}

\subsection*{You'll want to know: \ru{шить}}
\textbf{\ru{шить}} (\emph{shit}) — \textquotedblleft{}to sew\textquotedblright{}.

\begin{cousinweb}
\textbf{\ru{шить}} $\leftarrow$ Proto-Slavic \emph{*šiti} — inherited. Vasmer lists Lithuanian \emph{siūti}, Sanskrit \emph{sīvyati}, \textquotedblleft{}sews\textquotedblright{}, and Latin \emph{suō}, \textquotedblleft{}I sew\textquotedblright{}: English \emph{sew} is a cousin.
\end{cousinweb}

\begin{grammarlens}[title={What you've built}]
One more word for this chapter.
\end{grammarlens}

\subsection*{Your turn}
\begin{itemize}
\raggedright
  \item \emph{Say it:} \emph{shit}
  \item \emph{Say it:} \emph{shit}, once more
  \item \emph{Say it:} say \emph{shit}
  \item \emph{Recall:} say the Russian for a business trip, then the Russian for a contract, then say \emph{shit} again
\end{itemize}

\begin{tcolorbox}[breakable,colback=teal!4,colframe=teal!35!black,title={Before you move on}]
What does \ru{шить} mean? (\textquotedblleft{}To sew\textquotedblright{}.) Say it once more.
\end{tcolorbox}
\section[\texorpdfstring{\ru{резать} (rézat)}{резать (rézat)}]{\ru{резать} (rézat) — to cut}

\label{lesson:RU-C173-rezat}



\begin{quote}
Before the new one: say the Russian for a contract, then the Russian for to sew.
\end{quote}

\subsection*{You'll want to know: \ru{резать}}
\textbf{\ru{резать}} (\emph{rézat}) — \textquotedblleft{}to cut\textquotedblright{}.

\begin{cousinweb}
\textbf{\ru{резать}} $\leftarrow$ Proto-Slavic \emph{*rězati} — inherited. Vasmer relates it to Lithuanian \emph{rėžti}, \textquotedblleft{}to cut, to scratch\textquotedblright{}, and Greek \emph{rh\'{\={e}}gnymi}, \textquotedblleft{}I break\textquotedblright{}.
\end{cousinweb}

\begin{grammarlens}[title={What you've built}]
One more word for this chapter.
\end{grammarlens}

\subsection*{Your turn}
\begin{itemize}
\raggedright
  \item \emph{Say it:} \emph{rézat}
  \item \emph{Say it:} \emph{rézat}, once more
  \item \emph{Say it:} say \emph{rézat}
  \item \emph{Recall:} say the Russian for a contract, then the Russian for to sew, then say \emph{rézat} again
\end{itemize}

\begin{tcolorbox}[breakable,colback=teal!4,colframe=teal!35!black,title={Before you move on}]
What does \ru{резать} mean? (\textquotedblleft{}To cut\textquotedblright{}.) Say it once more.
\end{tcolorbox}
\section[\texorpdfstring{\ru{варить} (varít)}{варить (varít)}]{\ru{варить} (varít) — to boil, to cook}

\label{lesson:RU-C173-varit}



\begin{quote}
Before the new one: say the Russian for to sew, then the Russian for to cut.
\end{quote}

\subsection*{You'll want to know: \ru{варить}}
\textbf{\ru{варить}} (\emph{varít}) — \textquotedblleft{}to boil, to cook\textquotedblright{}.

\begin{cousinweb}
\textbf{\ru{варить}} $\leftarrow$ Proto-Slavic \emph{*variti}, from \emph{\ru{вар}}, \textquotedblleft{}boiling heat\textquotedblright{} — inherited. Vasmer relates \emph{\ru{вар}} to Lithuanian \emph{vìrti}, \textquotedblleft{}to boil\textquotedblright{}, and Latvian \emph{virt}, \textquotedblleft{}to boil\textquotedblright{}.
\end{cousinweb}

\begin{grammarlens}[title={What you've built}]
One more word for this chapter.
\end{grammarlens}

\subsection*{Your turn}
\begin{itemize}
\raggedright
  \item \emph{Say it:} \emph{varít}
  \item \emph{Say it:} \emph{varít}, once more
  \item \emph{Say it:} say \emph{varít}
  \item \emph{Recall:} say the Russian for to sew, then the Russian for to cut, then say \emph{varít} again
\end{itemize}

\begin{tcolorbox}[breakable,colback=teal!4,colframe=teal!35!black,title={Before you move on}]
What does \ru{варить} mean? (\textquotedblleft{}To boil, to cook\textquotedblright{}.) Say it once more.
\end{tcolorbox}
\section[\texorpdfstring{\ru{жарить} (zhárit)}{жарить (zhárit)}]{\ru{жарить} (zhárit) — to fry}

\label{lesson:RU-C173-zharit}



\begin{quote}
Before the new one: say the Russian for to cut, then the Russian for to boil.
\end{quote}

\subsection*{You'll want to know: \ru{жарить}}
\textbf{\ru{жарить}} (\emph{zhárit}) — \textquotedblleft{}to fry\textquotedblright{}.

\begin{cousinweb}
\textbf{\ru{жарить}} comes from \emph{\ru{жар}}, \textquotedblleft{}heat\textquotedblright{}, Proto-Slavic \emph{*žarŭ}, which Vasmer relates to Greek \emph{thermós}, \textquotedblleft{}warm\textquotedblright{}, and Sanskrit \emph{háras}, \textquotedblleft{}heat\textquotedblright{}.
\end{cousinweb}

\begin{grammarlens}[title={What you've built}]
One more word for this chapter.
\end{grammarlens}

\subsection*{Your turn}
\begin{itemize}
\raggedright
  \item \emph{Say it:} \emph{zhárit}
  \item \emph{Say it:} \emph{zhárit}, once more
  \item \emph{Say it:} say \emph{zhárit}
  \item \emph{Recall:} say the Russian for to cut, then the Russian for to boil, then say \emph{zhárit} again
\end{itemize}

\begin{tcolorbox}[breakable,colback=teal!4,colframe=teal!35!black,title={Before you move on}]
What does \ru{жарить} mean? (\textquotedblleft{}To fry\textquotedblright{}.) Say it once more.
\end{tcolorbox}
\section[\texorpdfstring{\ru{наливать} (nalivát)}{наливать (nalivát)}]{\ru{наливать} (nalivát) — to pour}

\label{lesson:RU-C173-nalivat}



\begin{quote}
Before the new one: say the Russian for to boil, then the Russian for to fry.
\end{quote}

\subsection*{You'll want to know: \ru{наливать}}
\textbf{\ru{наливать}} (\emph{nalivát}) — \textquotedblleft{}to pour\textquotedblright{}.

\begin{cousinweb}
\textbf{\ru{наливать}} is \emph{\ru{на}-}, \textquotedblleft{}into\textquotedblright{}, with \emph{-\ru{ливать}}, from \emph{\ru{лить}}, \textquotedblleft{}to pour\textquotedblright{}, which Vasmer relates to Lithuanian \emph{líeti}, \textquotedblleft{}to pour\textquotedblright{}.
\end{cousinweb}

\begin{grammarlens}[title={What you've built}]
One more word for this chapter.
\end{grammarlens}

\subsection*{Your turn}
\begin{itemize}
\raggedright
  \item \emph{Say it:} \emph{nalivát}
  \item \emph{Say it:} \emph{nalivát}, once more
  \item \emph{Say it:} say \emph{nalivát}
  \item \emph{Recall:} say the Russian for to boil, then the Russian for to fry, then say \emph{nalivát} again
\end{itemize}

\begin{tcolorbox}[breakable,colback=teal!4,colframe=teal!35!black,title={Before you move on}]
What does \ru{наливать} mean? (\textquotedblleft{}To pour\textquotedblright{}.) Say it once more.
\end{tcolorbox}
